\section{A localized change of variables in Hilbert space}
\label{sec:hilbert}

We now construct the Hilbert-space map required by
Proposition~\ref{prop:criterion}. The main issue is to let the output
coordinates separate input neighborhoods without losing global
invertibility.

\begin{theorem}
\label{thm:hilbert-map}
There are finite-dimensional Euclidean coordinate spaces $U_n$,
$U=(\bigoplus_{n\ge1}U_n)_2$, and a bijection
$h:U\oplus_2\ell_2\to U$ fixing zero, such that
\begin{equation}
\label{eq:h-bounds}
 \frac{24}{25}\|x-x'\|
 \le\|h(x)-h(x')\|
 \le\frac{26}{25}\|x-x'\|.
\end{equation}
Moreover, $D=h-\pi_U$ has local orthogonality: for every finite family
of distinct points $x_1,\ldots,x_m$ there are neighborhoods $O_j$ of
$x_j$, an integer $l$, and mutually orthogonal coordinate subspaces
$H_j\subset(\bigoplus_{n\ge l}U_n)_2$ such that
\[
 P_{\ge l}D(O_j)\subset H_j\qquad(1\le j\le m).
\]
Here $P_{\ge l}$ is the projection onto blocks with index at least
$l$, and a coordinate subspace is the closed span of a subset of the
specified orthonormal coordinates.
\end{theorem}

We begin with an onto isometric rearrangement determined by one unit
vector in each block. Allowing these vectors to depend on the input
will produce $h$. The rearrangement identifies the two requirements
on their motion: a weighted derivative bound will control metric
distortion, and separation of their coordinate supports will give
local orthogonality. We then construct the vectors, prove global
invertibility of finite approximants, and use an inverse-tail
estimate to show that their limiting map is onto.

\subsection{The frozen rearrangement}

For the moment let \(U=(\bigoplus_{n\ge1}U_n)_2\) have arbitrary
nonzero finite-dimensional real Hilbert blocks, let \(V=\ell_2\),
and put \(M=U\oplus_2V\). Choose a unit vector \(\xi_n\in U_n\)
for each \(n\). For \(x=(u,v)\in M\), write
\[
 \alpha_n=\langle \xi_n,u_n\rangle,\qquad
 u_n^\perp=u_n-\alpha_n\xi_n,\qquad
 C_{2m}=\alpha_m,\quad C_{2m-1}=v_m.
\]
Define
\begin{equation}
\label{eq:frozen-map}
 (L_\xi x)_n=u_n^\perp+C_n\xi_n.
\end{equation}
Thus the perpendicular component stays in block \(n\); the old slot
coefficient \(\alpha_m\) moves to the slot in block \(2m\), and
\(v_m\) fills the slot in block \(2m-1\). Figure~\ref{fig:slots}
shows this rearrangement.

\begin{figure}[ht]
\centering
\begin{tikzpicture}[
 >=Stealth,
 box/.style={draw=black!45,minimum width=40mm,minimum height=9mm,
             align=center,fill=black!2},
 every node/.style={font=\small}]
 \node[box] (perpin) at (0,1.7) {$u_n^\perp\in\xi_n^\perp$};
 \node[box] (alphain) at (0,0.35) {old slot $\alpha_m$};
 \node[box] (vin) at (0,-1) {extra coordinate $v_m$};
 \node[box] (perpout) at (7,1.7) {$u_n^\perp$ in block $n$};
 \node[box] (evenout) at (7,0.35) {$\alpha_m\xi_{2m}$ in block $2m$};
 \node[box] (oddout) at (7,-1) {$v_m\xi_{2m-1}$ in block $2m-1$};
 \draw[->,thick] (perpin)--node[above]{unchanged}(perpout);
 \draw[->,thick] (alphain)--node[above]{even slots}(evenout);
 \draw[->,thick] (vin)--node[above]{odd slots}(oddout);
 \node[font=\small\bfseries] at (0,2.5) {Input: $U\oplus_2\ell_2$};
 \node[font=\small\bfseries] at (7,2.5) {Output: $U$};
\end{tikzpicture}

\caption{The frozen map \(L_\xi\). The perpendicular components and
the two scalar sequences form orthogonal coordinates. The picture
describes the frozen linear isometry; when the slots depend on the
input, it gives the pointwise norm identity but not preservation of
distances between distinct inputs.}
\label{fig:slots}
\end{figure}

\begin{lemma}
\label[lemma]{lem:frozen-inverse}
\(L_\xi:M\to U\) is an onto linear isometry. If \(y=(y_n)\in U\)
and \(\beta_n=\langle\xi_n,y_n\rangle\), its inverse is
\begin{equation}
\label{eq:frozen-inverse}
 (L_\xi^{-1}y)_{U_m}
 =y_m-\beta_m\xi_m+\beta_{2m}\xi_m,\qquad
 (L_\xi^{-1}y)_{V_m}=\beta_{2m-1}.
\end{equation}
\end{lemma}

\begin{proof}
All formulas are linear for fixed \(\xi\). Since
\(u_n^\perp\perp\xi_n\), the squared norm of \eqref{eq:frozen-map}
is
\begin{align*}
 \|L_\xi x\|^2
 &=\sum_n\|u_n^\perp\|^2+\sum_n|C_n|^2\\
 &=\sum_n\bigl(\|u_n\|^2-|\alpha_n|^2\bigr)
       +\sum_m|\alpha_m|^2+\sum_m|v_m|^2
 =\|x\|_M^2.
\end{align*}
This also proves the output belongs to \(U\).

For \(y\in U\), \(|\beta_n|\le\|y_n\|\), so
\((\beta_n)\in\ell_2\). The first two terms
\(y_m-\beta_m\xi_m\) in \eqref{eq:frozen-inverse} are perpendicular
to its last term. The prescribed vector therefore belongs to \(M\),
since
\begin{align*}
 &\sum_m\left(\|y_m-\beta_m\xi_m+\beta_{2m}\xi_m\|^2
                   +|\beta_{2m-1}|^2\right)\\
 &\qquad=\sum_m\bigl(\|y_m\|^2-|\beta_m|^2\bigr)
                 +\sum_m|\beta_{2m}|^2+\sum_m|\beta_{2m-1}|^2
 =\|y\|^2.
\end{align*}
For this input the old slot coefficient is \(\beta_{2m}\), and the
perpendicular component in block \(m\) is \(y_m-\beta_m\xi_m\).
The even and odd output slots consequently recover every \(\beta_n\),
so \eqref{eq:frozen-map} returns \(y\). This proves surjectivity and
the inverse formula.
\end{proof}

\subsection{What the moving slots must satisfy}

Given continuously differentiable unit-vector maps $w_n:M\to U_n$,
define $h(x)=L_{w(x)}x$. Explicitly,
\begin{equation}
\label{eq:h-formula}
 \begin{split}
 h(x)_n&=u_n+w_n(x)\bigl(C_n(x)-\langle w_n(x),u_n\rangle\bigr),\\
 C_{2m}(x)&=\langle w_m(x),u_m\rangle,
 \qquad C_{2m-1}(x)=v_m.
 \end{split}
\end{equation}
The frozen isometry shows that this is well-defined and gives
\begin{equation}
\label{eq:pointwise-norm}
 \|h(x)\|=\|x\|.
\end{equation}
This identity does not yet compare distinct inputs, whose slots may
differ. To control that difference, fix $c=1/100$ and define the
nonnegative weights
\begin{equation}
\label{eq:weights-radii}
 a_n(x)^2=\|u_n\|^2+\|u_{\lceil n/2\rceil}\|^2
                         +|v_{\lceil n/2\rceil}|^2.
\end{equation}
These weights follow the routing in \eqref{eq:h-formula}:
$|C_n(x)|\le a_n(x)$ and $\|u_n\|\le a_n(x)$, so the scalar
coefficient multiplying $w_n(x)$ has absolute value at most $2a_n(x)$.
We shall choose the slots so that
\begin{equation}
\label{eq:slot-budget}
 \sum_{n\ge1}a_n(x)^2\|Dw_n(x)\|^2\le c^2
 \qquad(x\in M).
\end{equation}
We shall show that each derivative of a finite approximant is within
$4c$ of a frozen isometry.

The second requirement concerns the coordinate supports. For every
finite family of distinct input points, we shall find neighborhoods
whose slots use disjoint coordinates in all sufficiently high blocks.
Indeed, the $n$th block of $D(x)=h(x)-u$ is a scalar multiple of
$w_n(x)$, so this separation gives the required local orthogonality
of $D$. We now choose the block dimensions and construct slots
satisfying both requirements.

\subsection{Finite stages and separating coordinates}

Choose a nonnegative smooth function $\psi:\R\to\R$ that is
positive on $(-1,1)$ and zero outside $[-1,1]$. For $k\ge1$, put
\[
 \mathcal A_k=\{-k,-k+1,\ldots,k\}^k.
\]
For $a=(a_1,\ldots,a_k)\in\mathcal A_k$, define
\[
 G_{k,a}(t)=\prod_{j=1}^k\psi(kt_j-a_j)
 \qquad(t\in[-1,1]^k).
\]
Every point of $[-1,1]$ is within $1/(2k)$ of some $a_j/k$ with
$-k\le a_j\le k$. Thus $G_k(t)=(G_{k,a}(t))_{a\in\mathcal A_k}$
is never zero on the cube. Its Euclidean norm has a positive minimum
there. Each factor $t\mapsto\psi(kt-a_j)$ is supported on an interval
of length $2/k$.

Give $U_n$ an orthonormal basis consisting of a single stage-zero
vector and a separate set of coordinates indexed by $\mathcal A_k$
for each $1\le k\le n$. Different stages are orthogonal. In
particular, we may take
\begin{equation}
\label{eq:block-dimensions}
 \dim U_n=1+\sum_{k=1}^n(2k+1)^k.
\end{equation}
From now on, use these blocks in the separable Hilbert spaces
\[
 U=\left(\bigoplus_{n\ge1}U_n\right)_2,
 \qquad V=\ell_2,\qquad M=U\oplus_2V
\]
on which the slots will be constructed. Enumerate the scalar
coordinates of $M$ and apply $\tanh$ to each of them. This gives
smooth Lipschitz functions $p_j:M\to[-1,1]$ that separate points.

Let $W_{n,0}$ be the constant stage-zero unit vector. For $1\le k\le n$,
let $W_{n,k}(x)$ be the vector supported on stage $k$ with coordinates
\begin{equation}
\label{eq:normalized-stage}
 W_{n,k}(x)_a
 =\frac{G_{k,a}(p_1(x),\ldots,p_k(x))}
        {\|G_k(p_1(x),\ldots,p_k(x))\|_2}.
\end{equation}
It is a smooth unit-vector map. Compactness of the cube, the positive
lower bound for the denominator, and boundedness of the derivatives
of the $p_j$ give
\[
 \sup_x\|DW_{n,k}(x)\|\le L_k<\infty.
\]
The same $L_k$ works for every $n\ge k$, since their stage-$k$
coordinate formulas are identical. Put $L_0=0$. No bound uniform in
$k$ is needed.

\subsection{A common budget for motion}

To obtain \eqref{eq:slot-budget}, use the tails of the weights
\eqref{eq:weights-radii} as nonnegative radii:
\[
 r_n(x)^2=\sum_{j\ge n}a_j(x)^2.
\]
Every input coordinate in the second and third terms of
\eqref{eq:weights-radii} is counted twice, so
\begin{equation}
\label{eq:radii-identities}
 r_1(x)^2=3\|u\|^2+2\|v\|^2,
 \qquad a_n(x)^2=r_n(x)^2-r_{n+1}(x)^2.
\end{equation}
Each $r_n$ is the norm of a bounded linear map into a Hilbert space.
It is $\sqrt3$-Lipschitz, smooth where positive, and tends to zero
at every fixed $x$.

With the fixed constant $c=1/100$, set
\begin{equation}
\label{eq:gamma}
 \gamma(r)=
 \begin{cases}
 \displaystyle\frac{c}{r\log(1/r)},&0<r<e^{-2},\\[4pt]
 0,&r\ge e^{-2}.
 \end{cases}
\end{equation}
This function is nonincreasing as $r$ increases. It has an infinite
first integral near zero but a finite weighted square integral:
\begin{equation}
\label{eq:gamma-integrals}
 \int_0^b\gamma(r)\,dr=\infty\quad(b>0),
 \qquad
 \int_0^\infty2r\gamma(r)^2\,dr
 =2c^2\int_2^\infty t^{-2}\,dt=c^2.
\end{equation}
The infinite first integral permits arbitrarily many rotations; the
second estimate will bound their total effect on the derivative.

Choose successive compact intervals $[a_k',b_k']\subset(0,e^{-2})$
for the transitions from stage $k-1$ to stage $k$, with
\[
 b_{k+1}'<a_k',\qquad b_k'<2^{-k},\qquad
 \gamma(r)/2\ge L_{k-1}+L_k\quad(0<r\le b_k').
\]
They can be chosen together with smooth nonincreasing angles
$\theta_k:[0,\infty)\to[0,\pi/2]$ such that
\begin{equation}
\label{eq:angle-bounds}
 \theta_k(r)=\pi/2\ (r\le a_k'),\qquad
 \theta_k(r)=0\ (r\ge b_k'),\qquad
 |\theta_k'(r)|\le\gamma(r)/(2\sqrt3),
\end{equation}
and the angles are constant near the interval endpoints. Here is a
direct construction. First choose $b_k'$ small enough for the
derivative bound and the preceding interval. By
\eqref{eq:gamma-integrals}, a sufficiently small $a_k'>0$ makes the
integral of $\gamma/(2\sqrt3)$ over $(a_k',b_k')$ larger than
$\pi/2$. A smooth nonnegative function supported inside this
interval, bounded above by $\gamma/(2\sqrt3)$ and with integral
larger than $\pi/2$, exists by approximation from below. Rescale it
to integral $\pi/2$ and integrate from $r$ to $+\infty$ to obtain
$\theta_k$.

For each $n$, define a unit-vector map $P_n(x,r)$ by following these
same timings through stages $0,\ldots,n$ and then stopping. On the
$k$th transition interval, $k\le n$, put
\[
 P_n(x,r)=\cos\theta_k(r)\,W_{n,k-1}(x)
            +\sin\theta_k(r)\,W_{n,k}(x).
\]
Between intervals use the intervening single stage; use $W_{n,0}$
above $b_1'$ and $W_{n,n}$ below $a_n'$, including at $r=0$.
Orthogonality of the stages makes this a unit vector. The endpoint
choices make $P_n$ continuously differentiable. On a transition,
\eqref{eq:angle-bounds} and the bounds for $L_k$ give
\begin{equation}
\label{eq:P-derivatives}
 \|\partial_xP_n(x,r)\|\le\gamma(r)/2,
 \qquad
 \|\partial_rP_n(x,r)\|\le\gamma(r)/(2\sqrt3)
 \qquad(r>0).
\end{equation}
Between transitions the first bound follows from monotonicity of
$\gamma$, and the second derivative is zero. Stage zero is constant,
so both bounds also hold where $\gamma=0$.

Set
\begin{equation}
\label{eq:moving-slot}
 w_n(x)=P_n(x,r_n(x)).
\end{equation}
Although $r_n$ need not be differentiable where it vanishes, $w_n$
is continuously differentiable there: in a neighborhood of any such
point, $r_n<a_n'$ and $w_n=W_{n,n}$. At a point with $r_n>0$,
the chain rule and \eqref{eq:P-derivatives} give
$\|Dw_n(x)\|\le\gamma(r_n(x))$. These slots satisfy the required
budget \eqref{eq:slot-budget}. Indeed, when $r_n>0$, monotonicity gives
\[
 (r_n^2-r_{n+1}^2)\gamma(r_n)^2
 \le\int_{r_{n+1}}^{r_n}2r\gamma(r)^2\,dr,
\]
and summing uses \eqref{eq:gamma-integrals}. When $r_n=0$, also
$a_n=0$, so that term in \eqref{eq:slot-budget} is zero.

The derivative budget controls metric distortion. We still need the
independent locality property that will make the free-space
linearization completely continuous.

\subsection{Local orthogonality of the slots}

\begin{lemma}
\label[lemma]{lem:slot-locality}
For every finite family of distinct points of $M$ there are
neighborhoods $O_j$ and an integer $l$ such that the coordinate
supports of $\{w_n(x):x\in O_j,\ n\ge l\}$ are disjoint for
different $j$.
\end{lemma}

\begin{proof}
For an empty family, take $l=1$ and the empty family of neighborhoods.
For a singleton, take $l=1$ and $O_1=M$.
In either case there is no disjointness condition to check.
Now consider two distinct points. Some $p_j$ has different values at them.
Choose $k_0\ge j$ sufficiently large and neighborhoods on which all
values of that $p_j$ on the two sides are separated by more than
$2/k_0$. At any stage $k\ge k_0$, a coordinate label that occurs on
both sides would require the $j$th factor $\psi(kp_j-a_j)$ to
be nonzero at points $x,y$ on opposite sides. Then both $p_j(x)$
and $p_j(y)$ would lie in $((a_j-1)/k,(a_j+1)/k)$, contradicting
$|p_j(x)-p_j(y)|>2/k_0\ge2/k$. Thus the unions of the supports on
the two neighborhoods are disjoint at every common stage
$k\ge k_0$.

Fix $0<\rho<a_{k_0}'$, below the completed transition into stage
$k_0$. Choose $l\ge k_0$ so that $r_l$ at each center is less than
$\rho/2$, and shrink both neighborhoods to have radius less than
$\rho/(2\sqrt3)$. Since $r_n\le r_l$ for $n\ge l$ and $r_l$ is
$\sqrt3$-Lipschitz, we have $r_n(x)<\rho$ throughout either
neighborhood whenever $n\ge l$. Only stages $k\ge k_0$ are then
active in $w_n(x)$. Distinct stages are orthogonal, and supports at
the same stage are disjoint by the first paragraph. This proves the
two-point assertion. For finitely many points, intersect the
neighborhoods obtained for every pair and take the largest cutoff.
\end{proof}

The use of small neighborhoods here is essential: pointwise decay of
the $r_n$ and their common Lipschitz bound suffice. Uniform decay on
bounded subsets of $M$ is neither claimed nor needed.

The constructed slots now satisfy both the derivative budget and
local separation. Use them in \eqref{eq:h-formula} to define $h$.
It remains to prove the two-point bounds and surjectivity; the
pointwise norm identity alone does not give either conclusion.

\subsection{Finite approximants and global inverses}

Let $h_N$ be the same rearrangement with slots
\[
 w_n^{(N)}(x)=
 \begin{cases}w_n(x),&n\le N,\\W_{n,0},&n>N.\end{cases}
\]
Let $A_0$ denote the isometry with all slots equal to $W_{n,0}$.
Only output blocks $n\le2N$ of $h_N-A_0$ can be nonzero, and their
formulas are continuously differentiable. Thus $h_N$ is a
continuously differentiable map between Hilbert spaces. It also
satisfies $\|h_N(x)\|=\|x\|$.

\begin{lemma}
\label[lemma]{lem:approximant-inverse}
Each $h_N:M\to U$ is a continuously differentiable bijection with
continuously differentiable inverse, and
\begin{equation}
\label{eq:hN-bounds}
 (1-4c)\|x-x'\|\le\|h_N(x)-h_N(x')\|
             \le(1+4c)\|x-x'\|.
\end{equation}
\end{lemma}

\begin{proof}
Fix a point and a unit direction, and abbreviate the approximant slots
$w_n^{(N)}(x)$ and their directional derivatives by $w_n,\dot w_n$;
the slots with $n>N$ have derivative zero. Differentiating the inputs in
\eqref{eq:h-formula} with the slots frozen gives the onto isometry
$L_{w^{(N)}(x)}$. The remaining terms in output block $n$ are
\[
 \dot w_n(C_n-\langle w_n,u_n\rangle)
 +w_n\bigl(\dot C_n-\langle\dot w_n,u_n\rangle\bigr),
 \qquad
 \dot C_{2m}=\langle\dot w_m,u_m\rangle,
 \quad \dot C_{2m-1}=0.
\]
Put $R^2=\sum_n a_n(x)^2\|\dot w_n\|^2\le c^2$. Since
$|C_n-\langle w_n,u_n\rangle|\le2a_n$, the first term has
Hilbert norm at most $2R$. Each of the scalar sequences
$(\dot C_n)_n$ and $(\langle\dot w_n,u_n\rangle)_n$ has
$\ell_2$ norm at most $R$, because $a_n\ge\|u_n\|$.
Therefore
\begin{equation}
\label{eq:derivative-perturbation}
 \|Dh_N(x)-L_{w^{(N)}(x)}\|\le4c.
\end{equation}
A Neumann-series perturbation of the isometry gives
\begin{equation}
\label{eq:derivative-bounds}
 \|Dh_N(x)\|\le1+4c,
 \qquad \|Dh_N(x)^{-1}\|\le(1-4c)^{-1}.
\end{equation}

We next establish global bijectivity before using the lower
Lipschitz estimate. Put $F_N=\bigoplus_{n\le2N}U_n$ and let $P$
project onto $F_N^\perp$. We have
\[
 Ph_N=PA_0.
\]
For a fixed $b\in F_N^\perp$, the map $h_N$ thus restricts to
\begin{equation}
\label{eq:fiber-map}
 A_0^{-1}(b)+A_0^{-1}(F_N)\longrightarrow b+F_N.
\end{equation}
These affine spaces have the same finite dimension. The derivative
is an isomorphism between their tangent spaces. Indeed it maps
$A_0^{-1}(F_N)$ into $F_N$; conversely, for $z\in F_N$, the vector
$t=Dh_N(x)^{-1}z$ satisfies
$PA_0t=PDh_N(x)t=0$. Thus \eqref{eq:fiber-map} is a local
diffeomorphism. It is proper, since an inverse image of a compact
set is closed and norm-bounded by $\|h_N(x)\|=\|x\|$, hence compact
in this finite-dimensional source fiber.

A proper local diffeomorphism between affine Euclidean spaces of
the same dimension is bijective. To recall the topological argument,
its image is open by local invertibility and closed by properness,
so it is the entire target. Each fiber is compact and discrete,
hence finite. Local inverse neighborhoods, together with properness
to exclude further preimages nearby, make it a covering map. A
covering of a simply connected Euclidean space with connected
source has one sheet. Applying this to \eqref{eq:fiber-map} proves
that $h_N$ is bijective on each fiber and therefore on all of $M$.

The local inverses supplied by the inverse function theorem agree
and form a global continuously differentiable inverse. Integrating
the two bounds in \eqref{eq:derivative-bounds} along line segments,
for $h_N$ and for its now-defined inverse, proves
\eqref{eq:hN-bounds}.
\end{proof}

\subsection{Passing to an onto limit}

\begin{proof}[Completion of the proof of Theorem~\ref{thm:hilbert-map}]
For fixed \(x=(u,v)\), the first \(N\) output blocks of \(h_N(x)\)
and \(h(x)\) agree. Indeed, their own slots agree in those blocks,
and the old slot used in an even block \(2m\le N\) has \(m\le N\).
For either list \(\zeta=w(x)\) or \(\zeta=w^{(N)}(x)\), write
\[
 C_{2m}^{\zeta}=\langle\zeta_m,u_m\rangle,\qquad
 C_{2m-1}^{\zeta}=v_m.
\]
Then
\((L_\zeta x)_n-u_n
 =\zeta_n(C_n^\zeta-\langle\zeta_n,u_n\rangle)\).
Since \(|C_n^\zeta|\le a_n(x)\) and \(\|u_n\|\le a_n(x)\),
each difference from \(u_n\) has norm at most \(2a_n(x)\).
For the remaining blocks, the difference between the two maps has
norm at most \(4a_n(x)\). Hence
\begin{equation}
\label{eq:h-pointwise-limit}
 \|h_N(x)-h(x)\|
 \le4r_{N+1}(x)\longrightarrow0.
\end{equation}
For fixed \(x,x'\), take \(N\to\infty\) in
\eqref{eq:hN-bounds}. This proves the two-sided bounds
\eqref{eq:h-bounds}, since \(1-4c=24/25\) and
\(1+4c=26/25\). In particular, \(h\) is injective. It fixes zero by
\eqref{eq:pointwise-norm}.

We now prove surjectivity. Fix \(y=(y_n)\in U\) and, using
\Cref{lem:approximant-inverse}, put
\[
 x_N=h_N^{-1}(y).
\]
Freeze the slots of \(h_N\) at this particular input:
\[
 \xi_n^{(N)}=w_n^{(N)}(x_N),\qquad
 \beta_n^{(N)}=\langle \xi_n^{(N)},y_n\rangle.
\]
The directions may vary with \(N\). Nevertheless
\(y=L_{\xi^{(N)}}x_N\) is an exact equality for each \(N\), so
\Cref{lem:frozen-inverse} gives
\begin{equation}
\label{eq:xN-inverse}
 (x_N)_{U_m}
 =y_m-\beta_m^{(N)}\xi_m^{(N)}
       +\beta_{2m}^{(N)}\xi_m^{(N)},\qquad
 (x_N)_{V_m}=\beta_{2m-1}^{(N)}.
\end{equation}
The component \(y_m-\beta_m^{(N)}\xi_m^{(N)}\) is perpendicular
to \(\xi_m^{(N)}\). Thus its squared norm is
\(\|y_m\|^2-|\beta_m^{(N)}|^2\), and the last term in the
\(U_m\) formula adds orthogonally.

For every integer \(s\ge1\), it follows that
\begin{align}
\label{eq:inverse-tail}
 &\sum_{m\ge s}
   \left(\|(x_N)_{U_m}\|^2+|(x_N)_{V_m}|^2\right)\notag\\
 &\quad=\sum_{m\ge s}
   \left(\|y_m\|^2-|\beta_m^{(N)}|^2
               +|\beta_{2m}^{(N)}|^2
               +|\beta_{2m-1}^{(N)}|^2\right)\notag\\
 &\quad=\sum_{n\ge s}\|y_n\|^2
                -\sum_{n=s}^{2s-2}|\beta_n^{(N)}|^2
 \le\sum_{n\ge s}\|y_n\|^2.
\end{align}
Here \(|\beta_n^{(N)}|\le\|y_n\|\), so all squared-coefficient
series converge. The last equality is exact: the even and odd
indices \(2m,2m-1\) for \(m\ge s\) partition the integers
\(n\ge2s-1\); their sum cancels that tail of
\(\sum_{n\ge s}|\beta_n^{(N)}|^2\). The only uncancelled negative
terms have \(s\le n\le2s-2\). This finite sum is empty when \(s=1\).
Keeping it is what avoids double counting the output's perpendicular
and slot components.

The right side of \eqref{eq:inverse-tail} tends to zero with \(s\),
independently of \(N\). At \(s=1\) the equality gives
\(\|x_N\|=\|y\|\). For each \(s\), the heads
\[
 \bigl((x_N)_{U_m},(x_N)_{V_m}\bigr)_{m<s}
\]
lie in a bounded subset of the finite-dimensional space
\(\bigoplus_{m<s}(U_m\oplus\R)\). By successive subsequence
selection in these finite-dimensional spaces and a diagonal choice,
there is an increasing sequence \(N_j\) for which every fixed head
converges.

This diagonal subsequence is Cauchy in the full norm. Given
\(\epsilon>0\), choose \(s\) so that
\(\sum_{n\ge s}\|y_n\|^2<(\epsilon/4)^2\). By
\eqref{eq:inverse-tail}, the tail of each of \(x_{N_j}\) and
\(x_{N_k}\) has norm less than \(\epsilon/4\).
For sufficiently large \(j,k\), convergence of the fixed head makes
its difference smaller than \(\epsilon/2\). The triangle inequality
then gives \(\|x_{N_j}-x_{N_k}\|<\epsilon\).
Completeness of \(M\) yields \(x_{N_j}\to x\) in norm for some
\(x\in M\).

Finally, apply the pointwise convergence estimate at this fixed
limit \(x\), and use the common upper Lipschitz bound for \(h_{N_j}\):
\begin{align*}
 \|h(x)-y\|
 &\le\|h(x)-h_{N_j}(x)\|
      +\|h_{N_j}(x)-h_{N_j}(x_{N_j})\|\\
 &\le4r_{N_j+1}(x)+(1+4c)\|x-x_{N_j}\|\longrightarrow0.
\end{align*}
Thus \(h(x)=y\). Since \(y\) was arbitrary, \(h\) is onto.

The local orthogonality asserted in the theorem follows from
\Cref{lem:slot-locality}. For its neighborhoods $O_j$ and cutoff $l$,
let $H_j$ be the closed span of the union of coordinate supports of
$w_n(x)$ with $x\in O_j$ and $n\ge l$. These are mutually orthogonal
coordinate subspaces of $U_{\ge l}$. Each block $D(x)_n$ is a scalar
multiple of $w_n(x)$ by \eqref{eq:h-formula}. Its convergent tail
therefore belongs to $H_j$ for every $x\in O_j$, because $H_j$ is closed.
This verifies exactly \eqref{eq:local-orthogonality} and completes the proof.
\end{proof}
